\section{The modular at-most-two-stage theorem}
\label{sec:two-stage-composition}

\begin{theorem}[Containing-envelope composition]
\label{thm:two-stage-composition}
Fix $0<\alpha\leq1$, $0<\epsppr\leq1$, and one seed $s=e_v$.  Let a Stage~I
engine use charged work $W_{\rm disc}$ and return a containing-envelope
certificate $E$ for $\rho=\epsppr/2$.  Then
\cref{alg:two-stage-hybrid} returns $\widehat x$ satisfying
\cref{eq:two-stage-ppr-target} with total degree work
\begin{equation}
 W_{\rm total}
 \leq W_{\rm disc}
 +O\!\left(
   \frac{\vol(E)}{\sqrt\alpha}
   \log\frac{C_{E,v,\alpha}}{\epsppr}
 \right),
 \label{eq:two-stage-composition-work}
\end{equation}
where $C_{E,v,\alpha}$ is an explicit input/checkpoint scale and enters only
logarithmically.  In particular, if
$\vol(E)\leq C_E/\epsppr$, then
\begin{equation}
 W_{\rm total}
 =W_{\rm disc}
 +\widetilde O\!\left(
   \frac{C_E}{\epsppr\sqrt\alpha}
 \right).
 \label{eq:two-stage-product-tail}
\end{equation}
An exact-face certificate is the special case $E=A=S^\star$.
\end{theorem}

\begin{proof}
By \cref{prop:two-stage-envelope-linearization}, the exact restricted linear
solution $u_E$ already obeys
\begin{equation*}
 \norm{D^{-1/2}(x^0-\widetilde u_E)}_\infty\leq\rho.
\end{equation*}
The principal matrix $Q_E$ has spectrum in $[\alpha,1]$.  Accelerated
gradient or Chebyshev therefore reaches linear-system objective gap
$\alpha\tau^2/2$ in
\begin{equation*}
 O\!\left(\alpha^{-1/2}
   \log\frac{C_{E,v,\alpha}}\tau\right)
\end{equation*}
restricted matrix applications.  Strong convexity gives
$\norm{\widehat u_E-u_E}_2\leq\tau$.  Because $u_E\geq0$, taking the
coordinatewise positive part cannot increase this error.  Since every
$d_i\geq1$,
\begin{equation*}
 \norm{D_E^{-1/2}([\widehat u_E]_+-u_E)}_\infty\leq\tau.
\end{equation*}
Combining this with the envelope-linearization bound and
$\rho=\tau=\epsppr/2$ proves the semantic target.

One matrix application on $E$ costs $O(\vol(E))$.  For a point source the
initial restricted linear-system gap is bounded without a graph-size factor:
\begin{equation*}
 \frac12\langle b_E,Q_E^{-1}b_E\rangle
 \leq\frac{\norm{b_E}_2^2}{2\alpha}
 \leq\frac{\alpha}{2d_v}.
\end{equation*}
This supplies a valid logarithmic scale and completes the work bound.
\end{proof}

\begin{corollary}[Direct leakage-screen composition]
\label{cor:two-stage-leakage-composition}
Let Stage~I spend $W_{\rm screen}$ work and return a $\delta$ PPR
boundary-leakage certificate on $E$.  Refining the fixed principal solve to
normalized error $\tau$ returns a semantic PPR approximation whenever
$\delta+\tau\leq\epsppr$, with total work
\begin{equation}
 W_{\rm screen}
 +O\!\left(
   \frac{\vol(E)}{\sqrt\alpha}
   \log\frac{C_{E,v,\alpha}}{\tau}
 \right).
 \label{eq:two-stage-leakage-composition-work}
\end{equation}
This conclusion uses neither a regularization parameter nor containment of
$S^\star$.  In particular, if $E=V$ then the boundary gate is vacuous and
the statement reduces to one ordinary fixed solve.
\end{corollary}

\begin{proof}
Apply \cref{prop:two-stage-leakage-certificate} and the same fixed-system
convergence and triangle-inequality argument as in
\cref{thm:two-stage-composition}.
\end{proof}

\begin{theorem}[Screen-or-solve accounting]
\label{thm:two-stage-screen-or-solve}
Suppose Stage~I uses work $W_{\rm disc}$.  If it returns either
\begin{enumerate}[leftmargin=2.2em]
\item an exact RPPR point with $\rho\leq\epsppr$, or
\item a numerical gap/key certificate covered by
  \cref{cor:two-stage-numerical-handoff-returns} with
  $\rho+\delta\leq\epsppr$,
\end{enumerate}
then the algorithm returns after $W_{\rm disc}$ work, apart from writing the
sparse output.  Only a certificate that supplies a set without an equally
accurate numerical point pays the terminal term in
\cref{eq:two-stage-composition-work,eq:two-stage-three-budget-work}.
\end{theorem}

\begin{proof}
The exact RPPR point is $\rho$-accurate for PPR by
\cref{eq:two-stage-rppr-bridge}.  The second case is exactly
\cref{cor:two-stage-numerical-handoff-returns}.  The set-only case has no
claimed numerical output, so it invokes the containing- or
approximate-envelope composition theorem.
\end{proof}

\begin{corollary}[Approximate-envelope composition]
\label{cor:two-stage-three-budget}
Let Stage~I return a $\delta$-approximate-envelope certificate at
regularization $\rho$, and choose a terminal tolerance $\tau>0$ with
$\rho+\delta+\tau\leq\epsppr$.  Then the same linear Stage~II returns a
semantic $\epsppr$-PPR approximation, with terminal work
\begin{equation}
 O\!\left(
  \frac{\vol(E)}{\sqrt\alpha}
  \log\frac{C_{E,v,\alpha}}{\tau}
 \right).
 \label{eq:two-stage-three-budget-work}
\end{equation}
Thus a Stage~I engine may certify small exterior RPPR amplitude instead of
exact support containment.  Exact zero-key decisions are not part of the
minimal interface.
\end{corollary}

\begin{corollary}[An exact tail spends the full budget on screening]
\label{cor:two-stage-exact-tail-full-budget}
Suppose Stage~II solves its fixed principal system exactly, or to an error
accounted for outside the declared semantic tolerance.  Then the terminal
reserve may be zero.  A containing RPPR envelope may use
$\rho=\epsppr$, while a direct PPR truncation certificate may use
$\delta=\epsppr$.  In either case the returned point satisfies the semantic
target.  In the RPPR branch this also gives
\begin{equation*}
 \vol(S^\star(\epsppr))\leq\frac1{\epsppr}.
\end{equation*}
Thus an exact structural tail should replace the final Stage~I recovery
solve rather than be appended to it; the generic half-and-half split is
unnecessary for this lane.
\end{corollary}

\begin{proof}
The exact principal solution has no terminal numerical error.  The
envelope-linearization lemma bounds its normalized error by $\rho$, and the
direct leakage lemma bounds it by $\delta$, so the two claimed parameter
choices prove the semantic target directly.  The RPPR volume bound gives the
display.
The replacement statement is accounting rather than a new numerical step:
once the set certificate is available, only the ordinary principal system is
needed to produce the output.
\end{proof}

\begin{corollary}[An exact-face verifier may reuse its factorization]
\label{cor:two-stage-factor-reuse}
Suppose an exact-face verifier certifies $A=S^\star$ by factoring $Q_A$ and
solving $Q_Ap_A=c_{\rho,A}$.  The strict ordinary-PPR second stage solves
$Q_Au_A=b_A$ with the same factorization.  With a width-$w$ elimination
order, its incremental triangular-solve and output work is
$O(w^2\vol(A))$ after the $O(w^3\vol(A))$ factorization; for trees and
unicyclic graphs both are $O(\vol(A))$.  No obstacle or momentum state is
reused, only a certified factorization of the fixed linear operator.
\end{corollary}

\begin{proof}
The two systems have identical principal matrix and differ only in their
right-hand sides.  Reuse the stored elimination factors and perform forward
and backward substitution for $b_A$.  The stated sparse-factor costs follow
from the width bound.  Because the second right-hand side is the original PPR
load, the result is exactly the Stage~II point used in the composition
theorem.
\end{proof}

\begin{corollary}[The exact RPPR point is a certified linear warm start]
\label{cor:two-stage-rppr-warm-start}
Let $A=S^\star(\rho)$ and let $p_A=x_{\rho,A}^\star>0$.  For the ordinary
principal PPR system, the initial residual and correction are
\begin{equation}
 r_0=b_A-Q_Ap_A=\alpha\rho D_A^{1/2}\one,
 \qquad
 u_A-p_A=Q_A^{-1}r_0\geq0.
 \label{eq:two-stage-rppr-warm-residual}
\end{equation}
Moreover its linear-system objective gap is at most
\begin{equation*}
 \frac12 r_0^\top Q_A^{-1}r_0
 \leq\frac{\alpha\rho^2\vol(A)}2
 \leq\frac{\alpha\rho}{2}.
\end{equation*}
Thus a numerical terminal solve may start from $p_A$ with a fully observable
residual and a graph-size-free energy bound.  This is optional: when
$\rho=\epsppr$, $p_A$ is already a valid semantic output.
\end{corollary}

\begin{proof}
Every coordinate of the exact support is active, so stationarity gives
$Q_Ap_A=b_A-\alpha\rho D_A^{1/2}\one$.  This proves
\cref{eq:two-stage-rppr-warm-residual}; inverse positivity proves the sign.
Finally $Q_A\succeq\alpha I$ and
$\norm{D_A^{1/2}\one}_2^2=\vol(A)$ give the first inequality, while the RPPR
volume bound gives the second.
\end{proof}

\paragraph{For a genuinely set-only lane, the three budgets should be tuned,
not ritualized.}
The symmetric choice $\rho=\delta=\tau=\epsppr/3$ is transparent, but it is
not generally cost-minimizing.  For the exponential branch of the priority
monotone-SOR bound below, a useful leading-order proxy (suppressing logarithms
and fixed constants) is
\begin{equation}
 \frac{1}{\alpha\rho}+\frac{1}{\sqrt\alpha\,\delta},
 \qquad \rho+\delta=\epsppr-\tau.
 \label{eq:two-stage-budget-proxy}
\end{equation}
For a fixed terminal reserve $\tau$, this proxy is minimized by
\begin{equation}
 \frac{\rho}{\delta}=\alpha^{-1/4},\qquad
 \rho=\frac{\epsppr-\tau}{1+\alpha^{1/4}},\qquad
 \delta=\frac{\alpha^{1/4}(\epsppr-\tau)}{1+\alpha^{1/4}}.
 \label{eq:two-stage-tuned-budget}
\end{equation}
The minimized value is
$(\alpha^{-1/2}+\alpha^{-1/4})^2/(\epsppr-\tau)$.
Thus small $\alpha$ assigns most semantic error to RPPR regularization and a
smaller, but still explicit, amount to envelope omission.  The terminal
reserve affects an accelerated linear solve only logarithmically, so in a
 practical implementation it can be chosen from a modest fixed fraction or by
 the Stage~II residual target.  APPR has a different envelope law and should
 use its own split; the certificate race permits each lane to declare one.
 This is a worst-case proxy, not a claim that the tuned split dominates the
 symmetric choice instance by instance.  It must not be applied to a
 numerical Stage~I lane that can already invoke
 \cref{thm:two-stage-screen-or-solve}; for that lane the correct terminal
 reserve is zero and the output returns immediately.

\begin{corollary}[Exact support is not the main numerical bottleneck]
\label{cor:two-stage-support-not-needed}
Once any Stage~I construction returns an envelope of volume $O(1/\rho)$, all
remaining numerical work already has the target
$\widetilde O(1/(\rho\sqrt\alpha))$ scale.  Exact identification of the zero
coordinates and a second constrained optimization inside that envelope are
both unnecessary for semantic PPR.
\end{corollary}

\begin{corollary}[The obstacle history is disposable]
\label{cor:two-stage-history-disposable}
The terminal work bound requires only the set $E$ and the original PPR right
hand side $b_E$.  No Stage~I momentum, face-response bank, obstacle lower
checkpoint, strict-complementarity margin, Perron vector, or face-shock
transfer enters the theorem.  Such state may be used as a separately
certified warm start, but restarting the linear solve from zero already gives
\cref{eq:two-stage-composition-work}.
\end{corollary}
